\subsection{例：取整数值的多项式}

设 V 是有限维 \(\mathbf{Q}\)-向量空间，\(V^{*}\) 是其对偶，\(\mathcal{V}\) 是 V 中的一个格，\(\mathcal{V}^{*}\) 是对偶 \(\mathbf{Z}\)-模（其所对偶的格为 \(\mathcal{V}\)），它可典则地识别为 \(V^{*}\) 中的一个格，\(\mathbf{S}(V)\) 是 V 的对称代数，并且

\[
\lambda : \mathbf {S} (\mathrm{V}) \to \mathrm{A} (\mathrm{V} ^ {*})
\]

\(\mathbf{S}(\mathrm{V})\) 到 \(\mathrm{V}^*\) 上的多项式函数代数的典则双射（代数，第 IV 章，第 5 节，第 11 号，备注 1）。若将 \(\mathrm{A}(\mathrm{V}^* \times \mathrm{V}^*)\) 识别为 \(\mathrm{A}(\mathrm{V}^*) \otimes_{\mathbf{Q}} \mathrm{A}(\mathrm{V}^*)\)，则 \(\lambda\) 将 \(\mathbf{S}(\mathrm{V})\) 的余积变换为映射 \(\mathrm{A}(\mathrm{V}^*) \to \mathrm{A}(\mathrm{V}^* \times \mathrm{V}^*)\)，该映射把多项式函数 \(\varphi\)（定义在 \(\mathrm{V}^*\) 上）关联到多项式函数

\[
(x, y) \mapsto \varphi (x + y)
\]

在 \(\mathrm{V}^* \times \mathrm{V}^*\) 上（代数，第 IV 章，第 5 节，第 11 号，备注 2）。

记 \(\binom{\mathcal{V}}{\mathbf{Z}}\) 为 \(\mathbf{S}(\mathrm{V})\) 的子集，它由对应于从 \(\mathrm{V}^*\) 到 \(\mathbf{Q}\) 的多项式映射且在 \(\mathcal{V}^*\) 上取整数值的元素组成。

命题 2. (i) \(\begin{pmatrix} \mathcal{V} \\ \mathbf{Z} \end{pmatrix}\) 是双代数 \(\mathbf{S}(\mathrm{V})\) 中的双序，并且 \(\begin{pmatrix} \mathcal{V} \\ \mathbf{Z} \end{pmatrix} \cap \mathrm{V} = \mathcal{V}\)。

\begin{enumerate}
\def\labelenumi{(\roman{enumi})}
\setcounter{enumi}{1}
\item
  \(\mathbf{Z}\)-代数 \(\left( \begin{array}{c}\mathcal{V}\\ \mathbf{Z} \end{array} \right)\) 由 \(\left( \begin{array}{c}h\\ n \end{array} \right)\)（当 \(h\in \mathcal{V},n\in \mathbf{N}\) 时）生成。
\item
  若 \((h_1, \ldots, h_r)\) 是 \(\mathcal{V}\) 的一组基，则元素
\end{enumerate}

\[
\binom{h}{\mathbf {n}} = \binom{h _ {1}}{n _ {1}} \dots \binom{h _ {r}}{n _ {r}},
\]

其中 \(\mathbf{n} = (n_1, \ldots, n_r)\) 属于 \(\mathbf{N}^r\)，这些元素构成 \(\mathbf{Z}\)-模 \(\binom{\mathcal{V}}{\mathbf{Z}}\) 的一组基。

对 \(m \in \mathbf{N}\)，置 \(\mathbf{S}_{m}(\mathrm{V}) = \sum_{i \leq m} \mathbf{S}^{i}(\mathrm{V}), \mathbf{S}_{m}(\mathcal{V}) = \sum_{i \leq m} \mathbf{S}^{i}(\mathcal{V})\)。根据代数，第 IV 章，第 5 节，第 9 号，命题 15 及备注，

\[
\mathbf {S} _ {m} (\mathcal {V}) \subset \mathbf {S} _ {m} (\mathrm{V}) \cap \binom{\mathcal {V}}{\mathbf {Z}} \subset \frac {1}{m !} \mathbf {S} _ {m} (\mathcal {V})
\]

因此 \(\binom{\mathcal{V}}{\mathbf{Z}}\cap\mathrm{V}=\mathcal{V}\)。由于 \(\mathbf{S}_{m}(\mathcal{V})\) 是 \(\mathbf{S}_{m}(\mathrm{V})\) 中的一个格，\(\mathbf{S}_{m}(\mathrm{V})\cap\binom{\mathcal{V}}{\mathbf{Z}}\) 也是 \(\mathbf{S}_{m}(\mathrm{V})\) 中的一个格。另一方面，\(\mathbf{S}_{m}(\mathrm{V})\cap\binom{\mathcal{V}}{\mathbf{Z}}\) 是 \(\mathbf{S}_{m+1}(\mathrm{V})\cap\binom{\mathcal{V}}{\mathbf{Z}}\) 的直和因子（因为商无挠），因而有一个补空间，且该补空间是自由的 \(\mathbf{Z}\)-模。由此可知 \(\binom{\mathcal{V}}{\mathbf{Z}}\) 是自由的 \(\mathbf{Z}\)-模。显然，这是代数 \(\mathbf{S}(\mathrm{V})\) 中的一个幺序。设 \((u_{n})_{n\in\mathbf{N}}\) 是 \(\mathbf{Z}\)-模 \(\binom{\mathcal{V}}{\mathbf{Z}}\) 的一组基。它也是 \(\mathbf{Q}\)-模 \(\mathbf{S}(\mathrm{V})\) 的一组基，并且，对于所有

\[
\varphi \in \mathbf {S} (\mathrm{V} \times \mathrm{V}) = \mathbf {S} (\mathrm{V}) \otimes_ {\mathbf {Q}} \mathbf {S} (\mathrm{V}),
\]

存在唯一的序列 \((v_{n})\)，其元素属于 S(V)，使得 \(\varphi=\sum u_{n}\otimes v_{n}\)。如上，将 S(V) 识别为 A(V*)，并将 S(V) \(\otimes\) S(V) 识别为 A(V* × V*)。若 \(\varphi\in\binom{\mathcal{V}\times\mathcal{V}}{\mathbf{Z}}\)，则多项式函数 \(x\mapsto\varphi(x,y)\) 属于 \(\binom{\mathcal{V}}{\mathbf{Z}}\)，对所有 \(y\in\mathcal{V}^{*}\) 都成立。由此可知 \(v_{n}(y)\in\mathbf{Z}\) 对所有 n 及所有 \(y\in\mathcal{V}^{*}\) 都成立，换言之，\(v_{n}\in\binom{\mathcal{V}}{\mathbf{Z}}\)。这证明余积将 \(\binom{\mathcal{V}}{\mathbf{Z}}\) 映到 \(\binom{\mathcal{V}}{\mathbf{Z}}\otimes_{\mathbf{Z}}\binom{\mathcal{V}}{\mathbf{Z}}\)。若 \(h\in\mathcal{V}\) 且 \(n\in\mathbf{N}\)，则 \(\binom{h}{n}\) 将 \(u\in\mathcal{V}^{*}\) 映到整数 \(\binom{u(h)}{n}\)，所以 \(\binom{h}{n}\in\binom{\mathcal{V}}{\mathbf{Z}}\)。断言 (iii) 现在由将定理 1 应用于交换李代数 V 而得，(ii) 随之成立。

推论。令 \(\mathrm{X}\) 为一个不定元。多项式 \(\binom{\mathrm{X}}{n}\)，其中 \(n \in \mathbf{N}\)，构成以下 \(\mathbf{Z}\)-模的一组基：该模由 \(\mathrm{P} \in k[\mathrm{X}]\) 且满足 \(\mathrm{P}(\mathbf{Z}) \subset \mathbf{Z}\) 的多项式组成。

若 \(\mathrm{P}(\mathbf{Z}) \subset \mathbf{Z}\)，则拉格朗日插值公式（代数，第 IV 章，第 2 节，第 1 号，命题 6）表明 P 的系数属于 \(\mathbf{Q}\)；因此可以假设 \(k = \mathbf{Q}\)，并将命题 2 应用于 \(V = \mathbf{Q}\)、\(\mathcal{V} = \mathbf{Z}\)。

